\section{Introduction}
A panel starts a request, closes while the request is pending, and receives a result afterward. Whether the continuation should publish that result depends on the panel's requirements and on which operation still owns the output. The developer reviewing this code must answer a concrete question: which work becomes invalid at closure, and where must that invalidation constrain its continuation? Answering requires connecting the suspension point, lifecycle transition, publication, and protection.

Program comprehension during a change involves answering questions about relevant code and its relationships~\cite{sillito2006}. Temporal behavior adds a specific difficulty: a condition established before a boundary can cease to apply when execution resumes. Event-driven race detectors address related ordering problems using execution models~\cite{raychev2013}. A source review workflow can instead expose candidate relationships before the developer has reproduced an execution or established the application's validity requirements.

We organize this workflow around \emph{boundary review cards}. Each card links source references to four questions about a possible temporal relationship. A \emph{temporal contract} is an application requirement governing when an effect remains valid across execution or lifecycle boundaries. The analyzer selects candidate evidence; the reviewer establishes the requirement from repository documentation, implementations, or tests.

The tool provides three capabilities: (1) source extraction with explicit qualifications on candidate relations; (2) a four-question review interface that separates relation validity from protection sufficiency; and (3) shared, snapshot-based reports for command-line, editor, and agent-assisted review. The demonstration follows a pending request across panel closure and shows how the same review questions apply to an unguarded and a protected version.
